% =============================================================================
%  Chapter 6 — Implementation Summary and Design Decisions
% =============================================================================

\chapter{Implementation Summary and Design Decisions}
\label{ch:implementation}

This chapter collects the key design decisions made during the implementation
of simplex elements in \texttt{fall\_n}, with forward references to the
relevant source files and rationale grounded in the analysis of the preceding
chapters.


% ─────────────────────────────────────────────────────────────────────────
\section{Architecture Overview}
\label{sec:architecture}

\texttt{fall\_n} is a C++23 nonlinear finite element library.  Its element
hierarchy is fully \emph{compile-time polymorphic}: element types, quadrature
rules, constitutive models, and kinematic formulations are all resolved through
templates without virtual dispatch.

\begin{center}
\begin{tabular}{lll}
  \toprule
  Layer & Header & Responsibility \\
  \midrule
  Geometry        & \texttt{SimplexCell.hh}          & Reference coords, basis, topology \\
  Element         & \texttt{SimplexElement.hh}        & Iso-parametric map + integrator \\
  Integration     & \texttt{SimplexQuadrature.hh}     & Classical symmetric rules \\
                  & \texttt{StroudConicalProduct.hh}  & Arbitrary-order conical product \\
  Mesh I/O        & \texttt{GmshDomainBuilder.hh}     & Gmsh $\to$ \texttt{Domain} \\
  Model           & \texttt{Model.hh}                 & DOF management, BC application \\
  Post-processing & \texttt{VTKModelExporter.hh}      & Gauss $\to$ node projection, VTU \\
  \bottomrule
\end{tabular}
\end{center}


% ─────────────────────────────────────────────────────────────────────────
\section{Design Decisions}
\label{sec:decisions}


\subsection{Decision 1: HMS 4-Point as Default TET10 Quadrature}
\label{sec:dec_hms4}

\begin{description}[leftmargin=0pt,style=nextline]
  \item[Context] TET10 elements require at least degree-2 quadrature for the
    stiffness matrix.  The Keast 5-point rule (degree~3) was the initial
    default, following the common ``one degree above minimum'' heuristic.

  \item[Problem] The Keast 5-point rule has $w_1 = -2/15\,|\calS_3|$ at the
    centroid, causing the lumped $L^2$ projection to produce checkerboard
    artefacts (\cref{ch:negative_weights}).

  \item[Decision] Switch to HMS 4-point (degree~2) as the default.  Degree~2
    is the \emph{exact minimum} for the TET10 stiffness matrix; it provides
    no excess integration accuracy but avoids the negative-weight problem
    entirely.

  \item[Alternative considered] Conical product $n=2$ (8~points, degree~3,
    all positive).  This was rejected as the \emph{default} because it uses
    twice as many points as HMS 4, but it is available as
    \texttt{ConicalProductIntegrator<3,2>} for use cases requiring degree~3.

  \item[Source] \texttt{SimplexQuadrature.hh}, function
    \texttt{default\_simplex\_rule<3,2>()}.
\end{description}


\subsection{Decision 2: Golub--Welsch for Gauss--Jacobi Quadrature}
\label{sec:dec_gw}

\begin{description}[leftmargin=0pt,style=nextline]
  \item[Context] The Stroud conical product in 3D requires Gauss--Jacobi
    quadrature with parameters up to $\alpha = D-1 = 2$, $\beta = 0$.

  \item[Problem] An initial implementation used Newton iteration on the
    Jacobi polynomial (with Tricomi~/ asymptotic initial guesses).  For
    $\alpha = 2$ and $n \ge 3$, the iteration converged to duplicate roots
    because the initial guesses were insufficiently accurate for the
    distorted zero distribution of high-$\alpha$ Jacobi polynomials.

  \item[Decision] Replace Newton iteration entirely with the Golub--Welsch
    eigenvalue algorithm (\cref{sec:golub_welsch}).  This algorithm is
    unconditionally robust for all valid $\alpha, \beta > -1$ and any~$n$,
    since it requires no initial guesses.

  \item[Implementation] Uses Eigen's \texttt{SelfAdjointEigenSolver} on the
    $n \times n$ symmetric tridiagonal Jacobi matrix.  Cost is $O(n^2)$;
    since $n \le 10$ in practice, this is negligible.

  \item[Source] \texttt{StroudConicalProduct.hh}, function
    \texttt{gauss\_jacobi()}.
\end{description}


\subsection{Decision 3: Adaptive Projection as Default}
\label{sec:dec_spr}

\begin{description}[leftmargin=0pt,style=nextline]
  \item[Context] Gauss-point fields must be projected to nodes for VTU export.

  \item[Problem] Lumped $L^2$ projection is not robust on higher-order simplex
    elements.  The issue is deeper than negative quadrature weights:
    quadratic simplex vertex functions have signed support, so the row-sum
    lumped nodal measures can become non-positive even when \emph{all}
    quadrature weights are positive (\cref{ch:adaptive_projection}).

  \item[Decision] Make \texttt{compute\_material\_fields()} adaptive.  The
    exporter first inspects the local lumped nodal weights induced by the
    current element geometry and quadrature rule.  If every local weight is
    strictly positive, it uses lumped $L^2$ projection; otherwise it falls
    back to polynomial patch recovery
    (\texttt{patch\_recover\_to\_nodes()}).  The old volume-weighted averaging
    path is kept only as an explicit low-order fallback.

  \item[Trade-off] The adaptive pre-scan adds one post-processing pass over
    element quadrature points, but this cost is negligible compared with a
    solve and preserves the smoother lumped $L^2$ path for element families
    such as HEX27 where it is mathematically sound.  The patch-recovery path
    adds local least-squares work in post-processing only; it does not touch
    the FE assembly hot path.

  \item[Source] \texttt{VTKModelExporter.hh}, methods
    \texttt{lumped\_projection\_weights\_into()},
    \texttt{compute\_material\_fields()}, and
    \texttt{patch\_recover\_to\_nodes()}.
\end{description}


\subsection{Decision 4: Gmsh TET10 Node Reordering}
\label{sec:dec_reorder}

\begin{description}[leftmargin=0pt,style=nextline]
  \item[Context] Gmsh's local node ordering for 10-node tetrahedra differs
    from \texttt{fall\_n}'s ordering at the edge midpoints.

  \item[Problem] Without reordering, the Jacobian mapping is incorrect,
    producing negative determinants and garbage results.

  \item[Decision] Apply the permutation $(0,1,2,3,4,6,7,5,9,8)$ to map Gmsh
    indices to \texttt{fall\_n} indices.  This is done at mesh import time in
    \texttt{GmshDomainBuilder.hh}.

  \item[Verification] The simplex geometry tests
    (\texttt{test\_element\_geometry.cpp}) verify positive Jacobian
    determinants and correct isoparametric mapping for both TET4 and TET10
    elements imported from Gmsh.

  \item[Source] \texttt{GmshDomainBuilder.hh}, case 11 (TET10).
\end{description}


\subsection{Decision 5: PETSc Section Patch}
\label{sec:dec_petsc}

\begin{description}[leftmargin=0pt,style=nextline]
  \item[Context] \texttt{fall\_n} uses PETSc's DMPlex for distributed mesh
    management and \texttt{PetscSection} for DOF layout.

  \item[Problem] A bug in \texttt{PetscSectionSetConstraintIndices()} reads
    the atlas array with an incorrect index when \texttt{pStart} $> 0$
    (\cref{ch:petsc}).

  \item[Decision] Apply a local one-line patch:
    \texttt{s->atlasDof[point]} $\to$
    \texttt{s->atlasDof[point - s->pStart]}.

  \item[Impact] Without this patch, Dirichlet boundary conditions are silently
    incorrect on meshes where the vertex stratum does not start at point~0.

  \item[Source] PETSc~3.24.5, \texttt{src/sys/utils/section.c}, line~$\sim$2997.
\end{description}


\subsection{Decision 6: Hex27 VTK Face-Center Ordering}
\label{sec:dec_vtk_hex27}

\begin{description}[leftmargin=0pt,style=nextline]
  \item[Context] \texttt{fall\_n} stores the 27-node tensor-product HEX27 in
    lexicographic tensor order, while ParaView reads the exported VTU using
    \texttt{VTK\_TRIQUADRATIC\_HEXAHEDRON}.

  \item[Problem] The VTK exporter had the 6 face-center nodes partially
    permuted.  Corners and edge midpoints were correct, but ids 20--23 did not
    match VTK's actual \texttt{vtkTriQuadraticHexahedron} parametric-node
    ordering.  The visible symptom was a folded or diamond-like surface pattern
    on otherwise regular HEX27 meshes.

  \item[Decision] Correct the VTK permutation in
    \texttt{VTKCellTraits.hh}.  For the face centers, VTK expects the order
    $(x^-, x^+, y^-, y^+, z^-, z^+)$, which maps to the internal
    \texttt{fall\_n} indices $(12,14,10,16,4,22)$.

  \item[Verification] Add a regression test that compares the exporter
    permutation against the coordinates returned by
    \texttt{vtkTriQuadraticHexahedron::GetParametricCoords()}, instead of
    relying on a hand-written table alone.

  \item[Source] \texttt{src/post-processing/VTK/VTKCellTraits.hh},
    \texttt{tests/test\_simplex.cpp}.
\end{description}


% ─────────────────────────────────────────────────────────────────────────
\section{Nonlinear Showcase (\texttt{main.cpp})}
\label{sec:showcase}

The main executable exercises 7~constitutive/kinematic formulations on both
HEX27 and TET10 meshes (14~cases total) on a cantilever beam:

\begin{center}
\renewcommand{\arraystretch}{1.15}
\begin{tabular}{clllc}
  \toprule
  Case & Kinematics & Constitutive & Load & Steps \\
  \midrule
  1 & SmallStrain      & Linear elastic  & 0.05 & 1 (linear) \\
  2 & TotalLagrangian  & SVK             & 0.05 & 1 \\
  3 & TotalLagrangian  & SVK             & 0.50 & 5 \\
  4 & TotalLagrangian  & Neo-Hookean     & 0.50 & 5 \\
  5 & UpdatedLagrangian & SVK            & 0.05 & 1 \\
  6 & UpdatedLagrangian & Neo-Hookean    & 0.50 & 5 \\
  7 & SmallStrain      & J2 plasticity   & 1.00 & 20 \\
  \bottomrule
\end{tabular}
\end{center}

\paragraph{Material parameters:}
$E = 200$, $\nu = 0.3$, $\sigma_y = 0.250$, $H = 10.0$.

\paragraph{Geometry:}
$[0, 10] \times [0, 0.4] \times [0, 0.8]$.
Fixed face at $x = 0$, traction on face at $x = 10$.

\paragraph{Meshes:}
\begin{itemize}[nosep]
  \item \texttt{data/input/box.msh}: HEX27 (20~elements, 315~nodes).
  \item \texttt{data/input/box\_tet10.msh}: TET10 ($\sim$1329~elements,
        2840~nodes).
\end{itemize}

\paragraph{Post-processing:}
All cases use \texttt{compute\_material\_fields()} with adaptive dispatch
(lumped $L^2$ when safe, polynomial patch recovery otherwise), producing dual
VTU outputs (mesh + Gauss cloud) in
\texttt{data/output/}.


% ─────────────────────────────────────────────────────────────────────────
\section{Test Suite}
\label{sec:tests}

The simplex-related tests in \texttt{tests/test\_simplex.cpp} include 49~test
cases covering:
\begin{itemize}[nosep]
  \item Node count, reference coordinates, and shape function values for
        TET4 and TET10.
  \item Barycentric coordinate partition of unity and non-negativity.
  \item Shape function gradient consistency (finite difference vs.\ analytical).
  \item Face topology and local-to-global node mapping.
  \item Quadrature weight sums ($= 1/D!$ for the reference simplex).
  \item All-positive-weight verification for HMS and conical product rules.
  \item Points-inside-simplex verification for conical product.
  \item Polynomial exactness tests (degrees 3, 5, 7) for conical product
        integrators.
  \item \texttt{ConicalProductIntegrator<D, N>} template instantiation for
        $D \in \{1, 2, 3\}$, $N \in \{1, 2, 3, 4\}$.
\end{itemize}

All 49~tests and 16~element geometry tests pass as of the current revision.


% ─────────────────────────────────────────────────────────────────────────
\section{File Inventory}
\label{sec:file_inventory}

\begin{center}
\small
\renewcommand{\arraystretch}{1.1}
\begin{tabular}{p{6.8cm}p{7.5cm}}
  \toprule
  File & Purpose \\
  \midrule
  \texttt{src/geometry/SimplexCell.hh}
    & Reference simplex: coords, basis, topology (consteval) \\
  \texttt{src/elements/element\_geometry/SimplexElement.hh}
    & Element type combining geometry + integrator \\
  \texttt{src/numerics/.../SimplexQuadrature.hh}
    & Classical symmetric rules (HMS, Keast) \\
  \texttt{src/numerics/.../StroudConicalProduct.hh}
    & Golub--Welsch + Duffy + ConicalProductIntegrator \\
  \texttt{src/mesh/gmsh/GmshDomainBuilder.hh}
    & Gmsh import with TET10 node reordering \\
  \texttt{src/post-processing/VTK/VTKModelExporter.hh}
    & L2 and SPR projection, VTU export \\
  \texttt{data/input/box\_tet.geo}
    & Gmsh geometry for tetrahedral cantilever mesh \\
  \texttt{data/input/box\_tet10.msh}
    & TET10 mesh (2840 nodes) \\
  \texttt{tests/test\_simplex.cpp}
    & 49 simplex + conical product unit tests \\
  \texttt{main.cpp}
    & 14-case nonlinear showcase \\
  \bottomrule
\end{tabular}
\end{center}
